\section{Uncertainty, Confidence, and Risk}
\label{sec:entropy-confidence-risk}

This section treats uncertainty as the primary concept. Entropy is optional and
only meaningful when a discrete predictive distribution is defined.

Reliability diagnostics describe the trustworthiness of the current behavior
estimate. They are downstream from cost estimates and observations, but they are
not cost dimensions. Let
\[
\mathrm{Diag}(a)
=
h_{\mathrm{diag}}\!\left(
\widehat{\mathbf{C}}(a),
\{\epsilon_q(a)\},
z_a,
\theta_{k_a}
\right)
\]
denote a structured diagnostic object when the required inputs are available.
Diagnostics may include residual magnitude, missingness, sample scarcity, drift,
tail ratio, multimodality, calibration error, and estimator confidence. These
diagnostics belong to $D_{\mathrm{diag}}$ in the influence-domain system.

Risk is downstream from cost and diagnostics:
\[
\mathrm{Risk}(a)
=
g_{\mathrm{risk}}\!\left(
\widehat{\mathbf{C}}(a),
\mathrm{Diag}(a),
S(a);
\theta_{k_a}
\right),
\]
where $S(a)$ denotes consequence information such as tenant criticality,
dependency fragility, queue-starvation impact, or SLO sensitivity. Risk belongs
to $D_{\mathrm{risk}}$. It is not active resource demand, not worker occupancy,
and not merely low confidence. Low confidence can increase risk only when the
consequence of being wrong is operationally significant.

The separation between diagnostics and risk prevents a common modeling error. A
large residual says that the model and the observed workload disagreed. It does
not by itself say that the attempt is dangerous. A high-risk interpretation
requires consequence information: for example, whether underestimating this
attempt can starve a queue, overload a dependency, violate an SLO, or amplify a
retry storm.
